\documentclass{situationist-preprint}

\title{Aristotle on Slavery: A Test Case}
\msid{ETH-23-b}
\msfamily{ETH-23}
\msdiscipline{Normative and Applied Ethics}
\mstype{Consequence}
\mskeywords{Aristotle; natural slavery; moral anachronism; blame; epistemic distance; culpable ignorance}

\begin{document}
\maketitle

\begin{abstract}
Aristotle defended slavery, and his defense is the stock example of a great mind limited by its age. We apply the account of ETH-23-a, on which past acts are appraised by the correct standard and past agents are blameworthy only for what they could reasonably have known, and derive a definite verdict. Aristotle's defense of slavery and his participation in it were wrong. His epistemic distance from their wrongness was small: he reports the argument that slavery rests only on force, he concedes facts that undermine his own theory, and the capacities he denied to slaves were before him. The allowance for distance that he earns is therefore limited, and he is blameworthy to a substantial degree.
\end{abstract}

\section{The case}

In the first book of the \emph{Politics}, Aristotle argues that some human beings are slaves by nature: they lack the deliberative capacity and can share in reason only so far as to apprehend it in others, so that it is better for them to be ruled \citep[I.4--I.5, I.13]{AristotlePol}. He held slaves, and his will, as reported by Diogenes Laertius, provided for the freeing of several of them \citep[V.11--16]{DiogenesLaertius}. On the historicist view, Aristotle cannot be faulted for failing to see past the institutions of his world. On the presentist view, he was a defender of slavery and is to be blamed as one would blame such a defender today. ETH-23-a rejects both views (its Corollary~3.2), and Aristotle is the case it names as a test.

\section{Appraisal of the act}

By Principle~P1 of ETH-23-a, Aristotle's acts are appraised by the correct standard, of which present standards, insofar as they record moral progress, are our best approximation. Holding persons as property and defending the practice are wrong by that standard. By the first part of Theorem~3.1 of ETH-23-a, the verdict requires no allowance for epistemic distance.

\section{Appraisal of the agent}

\begin{proposition}\label{p:small}
Aristotle's epistemic distance from the wrongness of slavery was small.
\end{proposition}

\begin{proof}[Argument]
By Definition~2.3 of ETH-23-a, epistemic distance is the degree to which the evidence and arguments available to the agent failed to justify belief in the moral truth. Three facts show that they did not fail by much. First, the argument against slavery was available to him and stated by him: he reports that some held mastery to be contrary to nature, the distinction between slave and free to exist by convention only, and slavery therefore to be unjust, resting on force \citep[I.3]{AristotlePol}, and that many jurists challenged enslavement by conquest \citep[I.6]{AristotlePol}. Second, his theory required evidence that he concedes is lacking: he acknowledges that nature, though it intends to distinguish the bodies and souls of free and slave, often fails to do so, and that the souls are harder to see than the bodies \citep[I.5]{AristotlePol}. Third, the evidence of the capacities he denied was before him: he recommends that slaves be admonished rather than merely commanded \citep[I.13]{AristotlePol}, which presupposes the deliberative reason his theory withholds. A thinker of his acuity, attending to evidence he himself records, was in a position to conclude that the theory failed.
\end{proof}

\begin{proposition}\label{p:blame}
Aristotle is blameworthy for defending and practicing slavery to a substantial degree, reduced only modestly by his epistemic situation.
\end{proposition}

\begin{proof}
By Principle~P2 of ETH-23-a, an agent is blameworthy for a wrong act to the degree that she could reasonably have known it to be wrong. By Definition~2.2, what she could reasonably have known includes what inquiry she had reason to undertake would have shown. By Proposition~\ref{p:small}, the evidence and arguments available to Aristotle, including the counter-argument he reports, came close to justifying belief that slavery was wrong, and his own philosophical method gave him reason to examine the opposing view. The second part of Theorem~3.1 of ETH-23-a then yields a substantial degree of blame, with a modest reduction for the pervasive acceptance of the practice around him.
\end{proof}

The reduction is modest because general acceptance does not by itself alter the evidence. Where a practice is accepted because those who benefit from it decline to consider the arguments against it, the resulting ignorance is affected and does not excuse \citep{MoodyAdams1994}. Aristotle's case fits this description: the theory of natural slavery is an attempt to answer the objection he reports, and its defects lie in the very places where the evidence was available.

\begin{objection}
Aristotle's verdict was the product of a whole conceptual scheme, and the question whether he could have known better is idle; between his outlook and ours there is only a relativism of distance \citep{Williams1985}.
\end{objection}

\begin{reply}
Relativism of distance applies, if at all, where the alternatives were not real options for the agent. The view that slavery is unjust was a real option for Aristotle, since it was held by contemporaries and stated in his own text. His scheme did not exclude it; he considered it and rejected it with an argument whose premises he could see to be weak.
\end{reply}

\begin{objection}
On a sceptical view of culpable ignorance, an agent is blameworthy for acting on a false moral belief only if the belief resulted from an earlier culpable act, and such acts are hard to find \citep{Rosen2003}.
\end{objection}

\begin{reply}
Principle~P2 of ETH-23-a does not require an earlier act of akrasia, only that the agent could reasonably have known the act to be wrong. Aristotle's ignorance, if it was ignorance, persisted in the face of arguments and evidence he possessed and considered. That suffices.
\end{reply}

\section{Conclusion}

Aristotle acted wrongly in defending and practicing slavery, and he is substantially blameworthy for it. The test case confirms the implication of ETH-23-a that the allowance for epistemic distance is often smaller than historicists suppose: the arguments that later prevailed were, in his case, available and known to him.

\begin{thebibliography}{9}\small
\bibitem[Aristotle(Pol.)]{AristotlePol} Aristotle. \emph{Politics}. Cited by book and chapter.
\bibitem[Diogenes Laertius(Lives)]{DiogenesLaertius} Diogenes Laertius. \emph{Lives of Eminent Philosophers}. Cited by book and section.
\bibitem[Moody-Adams(1994)]{MoodyAdams1994} Moody-Adams, M.~M. (1994). Culture, responsibility, and affected ignorance. \emph{Ethics}, 104(2).
\bibitem[Rosen(2003)]{Rosen2003} Rosen, G. (2003). Culpability and ignorance. \emph{Proceedings of the Aristotelian Society}, 103.
\bibitem[Williams(1985)]{Williams1985} Williams, B. (1985). \emph{Ethics and the Limits of Philosophy}. Fontana.
\end{thebibliography}

\end{document}
